\section{Global identification of marked contact germs}
\label{sec:v66-global-contact}

The local analytic coordinate theorem of the preceding section controls
perturbations of complete germs.  Exact identification has a stronger
form: for a fixed marked polygon there is no second positive-curvature
branch with the same actions.  The point is already visible at degree
two.  A Schur-complement identity converts the curvature inverse into a
strict contraction whose constant is determined by the observed action
Hessians.  The signed higher-degree recursion then gives exact uniqueness
without a closeness assumption.  The analytic estimates remain local.

Throughout this section the polygon, graph coordinates and signs are
those of Section~\ref{sec:v65-periodic-inverse}.  In particular
$k_i=L_i^{-1}>0$, $c_i=F_i''(0)>0$ and $s_i=(S_i^-)''(0)$;
indices are modulo the physical contact period $r$.  We regard $k$ as
fixed and write $\mathcal R_k(c)=s$.  This is a map of quadratic data,
independent of the higher contact jets.  Its definition uses the actual
half-line construction, not an arbitrary solution of a Riccati equation.

\subsection{The quadratic inverse}

\begin{lemma}[Boundary Schur complement]\label{lem:v66-riccati}
For every positive curvature vector $c$, the actual half-line action
Hessians satisfy $0<c_i<s_i$ and
\begin{equation}\label{eq:v66-schur}
 s_i=k_i+c_i-\frac{k_i^2}{k_i+c_{i+1}+s_{i+1}},\qquad
 a_i=\frac{k_i}{k_i+c_{i+1}+s_{i+1}}\in(0,1).
\end{equation}
Here $a_i$ is the unsigned one-step contraction in
Lemma~\ref{lem:v65-signed-jacobi}.  Conversely, positive vectors $c,s$
satisfying the first identity have $s=\mathcal R_k(c)$.
\end{lemma}
\begin{proof}
The half-line starting at $i$ is obtained by adjoining its first edge
to the half-line starting at $i+1$.  At quadratic order its stationary
value is the minimum over $y$ of
\[
 \tfrac12(k_i+c_i)u^2-\epsilon_i k_iuy
                  +\tfrac12(k_i+c_{i+1}+s_{i+1})y^2.
\]
The strictly positive half-line Hessian established in the preceding
section gives $s_{i+1}>0$.  Eliminating $y$ proves the first identity;
the stationary ratio is $y/u=\epsilon_i a_i$.  The denominator exceeds
$k_i$, proving both $a_i<1$ and $s_i>c_i$.

For the converse, define $a_i$ by the second identity.  Substituting the
first identity at $i+1$ gives
\[
 \frac{k_i}{a_i}
   =k_i+k_{i+1}+2c_{i+1}-k_{i+1}a_{i+1}.
\]
Thus the products $y_j=a_i\cdots a_{i+j-1}u$ solve the sign-conjugated
Jacobi equation with boundary value $u$.  They decay exponentially
because there are finitely many phases and every $a_i<1$.  Coercivity
makes this decaying solution unique: the difference of two such solutions
has zero boundary value, and summation of its positive Jacobi quadratic
form gives zero.  Restoring the signs gives the actual linearized
half-line.  Its initial action Hessian is $k_i+c_i-k_i a_i=s_i$.
Consequently these are the actual action Hessians.
\end{proof}

For $s\in(0,\infty)^r$ and $z\in[0,\infty)^r$, put
\begin{equation}\label{eq:v66-projection-map}
 \begin{aligned}
 (\Phi_s z)_i&=
 \left[s_i-k_i+\frac{k_i^2}{k_i+s_{i+1}+z_{i+1}}\right]_+,
 \qquad [t]_+=\max(t,0),\\
 q(s)&=\max_i\left(\frac{k_i}{k_i+s_{i+1}}\right)^2<1.
 \end{aligned}
\end{equation}
The positive-part operation specifies an inverse algorithm even for
positive data that are not realizable by strictly positive curvatures.
It does not add zero-curvature tables to the geometric class.

\begin{theorem}[Global curvature coordinates]\label{thm:v66-curvature}
The map $\mathcal R_k:(0,\infty)^r\longrightarrow(0,\infty)^r$
is one-to-one and a real analytic diffeomorphism onto its open image.
That image and the inverse admit the following description.

For every $s>0$, the map $\Phi_s$ has a unique fixed point $z(s)$ in
the closed positive orthant.  Its iterates from any $z^{(0)}\ge0$
converge to $z(s)$, with
\begin{equation}\label{eq:v66-iteration-error}
 \|z^{(m)}-z(s)\|_\infty
 \le \frac{q(s)^m}{1-q(s)}
                    \|z^{(1)}-z^{(0)}\|_\infty\quad(m\ge0).
\end{equation}
The first-increment estimate \eqref{eq:v66-iteration-error} is an
\emph{a priori} bound.  A residual at the current iterate gives the
\emph{a posteriori} bound
\begin{equation}\label{eq:v67-residual-error}
 \|z^{(m)}-z(s)\|_\infty\le E_m(s),\qquad
 E_m(s):=\frac{\|\Phi_s z^{(m)}-z^{(m)}\|_\infty}{1-q(s)}.
\end{equation}
For exact iteration its numerator is
$\|z^{(m+1)}-z^{(m)}\|_\infty$.
Precisely,
\begin{equation}\label{eq:v66-image}
 s\in\mathcal R_k((0,\infty)^r)
       \quad\Longleftrightarrow\quad z_i(s)>0\text{ for every }i;
 \qquad \mathcal R_k^{-1}(s)=z(s).
\end{equation}
For any $s,\widetilde s>0$, define the observable constant
\[
 q_* = \max_i
 \left(\frac{k_i}{k_i+\min(s_{i+1},\widetilde s_{i+1})}\right)^2<1.
\]
Then, including for data outside the geometric image,
\begin{equation}\label{eq:v66-pair-lipschitz}
 \|z(s)-z(\widetilde s)\|_\infty
       \le\frac{1+q_*}{1-q_*}\|s-\widetilde s\|_\infty.
\end{equation}
\end{theorem}
\begin{proof}
For $z\ge0$ the expression inside the positive part is strictly smaller
than $s_i$.  Hence $\Phi_s$ maps the closed orthant into
$\prod_i[0,s_i]$.  Its untruncated expression has derivative in
$z_{i+1}$ equal to
$-k_i^2/(k_i+s_{i+1}+z_{i+1})^2$, of absolute value at most $q(s)$.
The positive part is one-Lipschitz.  The map is therefore a strict
contraction in the maximum norm on the complete closed orthant.
The contraction theorem and summation of the successive differences
prove existence, uniqueness and \eqref{eq:v66-iteration-error}.
For any $v\ge0$, the triangle inequality and contraction give
\[
 \|v-z(s)\|_\infty
 \le\|v-\Phi_s v\|_\infty+q(s)\|v-z(s)\|_\infty.
\]
This proves \eqref{eq:v67-residual-error}.  It also provides a bound
for an inexact evaluation: if $w$ satisfies
$\|w-\Phi_s v\|_\infty\le\eta$, replace the residual numerator
by $\|w-v\|_\infty+\eta$.  An error bound for evaluating the map,
when used, is part of this certificate and is not assumed to be zero.

If $s=\mathcal R_k(c)$ for $c>0$, Lemma~\ref{lem:v66-riccati} says
that $c=\Phi_s c$.  Therefore $c=z(s)$, which proves global
injectivity.  If $z(s)>0$, no positive part is active at the fixed
point, so \eqref{eq:v66-schur} holds with $c=z(s)$.  The converse
in that lemma proves the other implication in \eqref{eq:v66-image}.
This proves the image characterization without identifying formal
positive data with geometric data.

For the two-point estimate, changing $s_i$ contributes at most
$\|s-\widetilde s\|_\infty$ to $\Phi_s z$; changing $s_{i+1}$
in its denominator contributes at most
$q_*\|s-\widetilde s\|_\infty$.  All intermediate denominators
are bounded below by the one used in $q_*$.  Changing $z$ contributes
at most $q_*\|z-\widetilde z\|_\infty$.  Apply these bounds at
the two fixed points and move the last term to the left.  This proves
\eqref{eq:v66-pair-lipschitz}.

On the image the fixed point is strictly positive.  Locally the positive
part can be omitted, and the remaining equation is real analytic in
$(s,c)$.  Its derivative in $c$ is $I+T_2$, where
$(T_2 v)_i=a_i^2v_{i+1}$, and has a convergent Neumann inverse because
$\|T_2\|_\infty<1$.  Thus the inverse is locally real analytic.
The forward half-line Hessian is real analytic in $c$: for example the
positive half-line Jacobi operator with periodic coefficients has a locally
uniformly convergent
resolvent expansion, and the boundary Schur complement is its analytic
matrix element.  Differentiating \eqref{eq:v66-schur} also gives
\begin{equation}\label{eq:v66-differential}
 (I-T_2)\,ds=(I+T_2)\,dc.
\end{equation}
Both factors are invertible.  Hence the forward map is a local real
analytic diffeomorphism at every $c>0$.  Its image is open, and the
already proved injectivity makes these local inverses one global inverse
on that image.
\end{proof}

The projection in \eqref{eq:v66-projection-map} is useful only in the
reconstruction algorithm; it plays no role in the analytic inverse near
a realized datum.  Formula~\eqref{eq:v66-iteration-error} also gives a
finite sufficient test for membership in the image: if every component
of an iterate exceeds either certified error bound, then $z(s)>0$.
For the residual bound this test is $\min_i z_i^{(m)}>E_m(s)$.
It certifies membership in the quadratic image, not extension to a
prescribed collection of globally disjoint analytic obstacles.
The constants are allowed to deteriorate as an action Hessian tends to
zero.  No uniform assertion at a grazing or zero-curvature limit is used.

There is a second direct uniqueness check.  If two positive curvature
vectors have the same $s$, subtraction of \eqref{eq:v66-schur} gives
\[
 c_i-\widetilde c_i
       =-a_i\widetilde a_i(c_{i+1}-\widetilde c_{i+1}).
\]
Iteration around the cycle forces the difference to vanish, since the
product of the $a_i\widetilde a_i$ is strictly smaller than one.
This argument treats odd and even physical periods alike.  It uses only
quadratic data; the signs in odd higher contact orders still have to be
retained in the geometric reconstruction.

For a normal two-contact channel, $k_0=k_1=1/g$ and $c_i$ are the
physical curvatures.  Set $H=(1+g^2s_0s_1)^{1/2}$.  The recovered
curvatures are
\[
 c_0=\frac{H\sqrt{s_0/s_1}-1}{g},\qquad
 c_1=\frac{H\sqrt{s_1/s_0}-1}{g}.
\]
These are precisely the alternating curvature formulas of
Theorem~\ref{thm:v26-density-inverse}.  Indeed put
$u_i=gs_i$ and $C_i=1+gc_i$; then $C_0C_1=1+u_0u_1$ and
$C_0/C_1=u_0/u_1$, which give
$(C_0-u_0)(C_1+u_1)=1$ and its reversed version.
These are the two identities in \eqref{eq:v66-schur}.
Whenever the displayed curvatures are positive, uniqueness identifies
them with the fixed point.  This is a specialization of the geometric
inverse, not a reduction of the two observation models.

Not every positive vector is an action datum for positive curvatures.
For example $r=2$, $k_0=k_1=1$ and $s=(1/10,10)$ give
$z(s)=(0,109/11)$.  The first component of the untruncated map is
negative for every $z\ge0$, and the second fixed-point equation gives
the stated value.  Thus this $s$ is excluded by \eqref{eq:v66-image},
even though the stable reconstruction algorithm is defined there.

\subsection{Exact rigidity without a local branch choice}

\begin{theorem}[Global marked contact identification]
\label{thm:v66-global-rigidity}
Fix a clear nongrazing marked polygon with distinct physical contact
points as in Section~\ref{sec:v65-periodic-inverse}.  For every $M\ge2$
the map from positive-curvature smooth contact jets through order $M$
to the future action jets through that order is one-to-one, with a smooth
inverse on its image.  No closeness between the candidate jets is required.

Consequently, the phase-resolved two-offset limiting laws
\eqref{eq:v65-two-laws}, on their common positive coordinate square,
identify every finite smooth contact jet.  Among positive-curvature
analytic contact tuples with this same polygon they identify the complete
contact germs, without a local-neighborhood hypothesis.  The two offsets
share their endpoint amplitude factors as in that formula; the amplitudes,
curvatures and graph heights are not supplied.

If two actual tables have the same marked polygon, fixed lattice and
obstacle labels, their boundaries are connected closed embedded analytic
curves, and the polygon visits every labelled obstacle, equality of these
laws implies equality of the complete labelled obstacle images.  The
tables need not be close.  The fixed placement and full visitation are
part of this conclusion.
\end{theorem}
\begin{proof}
At degree two apply Theorem~\ref{thm:v66-curvature}.  This determines
$c$, hence the actual $a_i$ by \eqref{eq:v66-schur}, and the signed
factors $\sigma_i=\epsilon_i a_i$.  At degree $n\ge3$ use the
actual-smooth recursion from Theorem~\ref{thm:v65-cyclic-jets}:
\[
 q_n=(I-T_n)(I+T_n)^{-1}
                 \{s_n-\mathcal R_n(q_2,\ldots,q_{n-1})\}.
\]
Here $T_n$ is the signed shift of \eqref{eq:v65-shift}, now fully
known from the reconstructed quadratic jets and the supplied polygon.
Its inverse exists at every degree.  Induction determines all higher
jets uniquely.  The lower-degree remainders are functions of actual
smooth jets by Lemma~\ref{lem:v65-envelope}, so this is not an
identification of unrelated formal expansions.  The curvature inverse is
smooth on its image, and the remaining finite recursion is smooth, proving
the finite-order assertion.  Arbitrary finite higher jets with $c>0$
have local smooth graph representatives; after restricting their collars
they satisfy the same local geometric conditions.

Proposition~\ref{prop:v65-two-offset} extracts the physical actions
$A_i,C_i$ from the two laws.  The known momenta give
$S_i^-=A_i+p_i x$.  Thus equality of laws gives equality of every
future action jet and hence of every contact jet.  In the analytic class
the Taylor series of each contact graph converges on some neighborhood
of zero.  Equality of all its coefficients gives equality of the graph
germs.  One may restrict two different analytic radii to a smaller common
one; no comparison in a prescribed local inverse ball is needed for this
exact assertion.  In the smooth category we conclude equality of jets,
not equality of germs from flatness alone.

For actual analytic tables this supplies a common open arc on every
visited obstacle.  The continuation argument in
Corollary~\ref{cor:v65-law-rigidity} then identifies each entire
connected closed boundary image.  More explicitly, at a limit point of
an interval of common arc both curves have the same tangent.  They are
analytic graphs over that tangent on a common neighborhood, and their
identity continues across the limit point.  Compactness and connectedness
continue it around the boundary.  All obstacles are visited and both
lattice and placement were supplied, so the labelled tables coincide.
\end{proof}

Theorem~\ref{thm:v66-global-rigidity} strengthens the exact uniqueness
clauses in Corollary~\ref{cor:v65-law-rigidity}; the latter remains
useful for its physical derivation and local estimates.  The new exact
argument does not replace the complete function-space inverse in
Theorem~\ref{thm:v65-analytic-inverse}.  That theorem controls all
lower-degree couplings in a fixed analytic norm, which coefficient
uniqueness alone cannot do.  In particular the analytic prior and radius
loss in Theorem~\ref{thm:v65-real-stability} are unchanged.

\begin{proposition}[Finite-flight reconstruction on a positive class]
\label{prop:v66-finite-flight}
Fix $M\ge2$ and a compact class of marked smooth contact data with
positive curvature, the fixed-order smooth bounds of the relative law,
and the positive-density, residual and separated-offset margins of
Proposition~\ref{prop:v65-two-offset}.  Use the same fixed polygon.
For sufficiently small errors, the two-offset densities of $N$-flight
bridges, with $N$ a multiple of a returning oriented period, recover the
contact jets through order $M$ with error at most
\begin{equation}\label{eq:v66-fixed-order-error}
 C_M(\tau^N+\varepsilon).
\end{equation}
Here $\varepsilon$ is the maximum interior $C^M$ error at the two
offsets and every phase.  The quadratic reconstruction is the fixed
point \eqref{eq:v66-projection-map}, not a selected local solution of
the curvature equation.  All later degrees use the signed recursion.

More precisely, write $t=(t_2,\ldots,t_M)$ for the extracted future
action jets and stop the quadratic iteration at $v=z^{(m)}(t_2)$.
Use the off-fixed-point recursion specified in the proof.  For
sufficiently small $\tau^N+\varepsilon+E_m(t_2)$ it is defined and
satisfies, in the maximum norm of the derivative coefficients,
\begin{equation}\label{eq:v67-complete-stopping-error}
 \|\widehat q_{\le M}^{(m)}-q_{\le M}\|_\infty
 \le C_M(\tau^N+\varepsilon)+L_M E_m(t_2)
 \le\widetilde C_M\{\tau^N+\varepsilon+E_m(t_2)\}.
\end{equation}
The constants are uniform on the stipulated compact class, with $M$
fixed.  The first-increment bound in \eqref{eq:v66-iteration-error}
may replace $E_m$.  The factor $L_M$ includes the sensitivity of
all the lower-degree couplings in the higher-jet reconstruction.
\end{proposition}
\begin{proof}
The relative physical law first compares the finite-flight densities
with their realizable limits, with $C^M$ error $O(\tau^N)$ uniformly
on the compact class.  Stable rational extraction then gives the gauged
future action jets with error $O(\tau^N+\varepsilon)$.  Subtract their
constant and linear Taylor terms, as in the preceding section; the true
gauged action is fixed by this bounded projection.  The true quadratic
data have a common positive floor.  For sufficiently small errors the
estimated $s$ remain positive, so \eqref{eq:v66-pair-lipschitz} applies
with a uniform $q_*<1$.  Its output is close to the true $c$, whose
components also have a positive floor.  The reconstructed $c$ are
therefore positive, and the projection is inactive at their fixed point.

At each of the finitely many remaining orders the coefficients, the
inverse signed block and the lower-jet remainder are smooth on bounded
sets of jets with a positive curvature margin.  Inductively their error
is $O(\tau^N+\varepsilon)$.  The finite constants can be bounded on
a common compact jet box; the order $M$ is fixed throughout.  This proves
\eqref{eq:v66-fixed-order-error}.

We specify the finite-iteration reconstruction, since the error in
curvature must pass through the subsequent nonlinear steps.  Keep the
extracted action vector $t$ fixed.  For $v>0$ near $z(t_2)$ set
\begin{equation}\label{eq:v67-off-fixed-block}
 \begin{split}
 \widehat\sigma_i(v;t_2)&=
      \frac{\epsilon_i k_i}{k_i+t_{2,i+1}+v_{i+1}},\\
 (\widehat T_n(v;t_2)w)_i&=
                         \widehat\sigma_i(v;t_2)^n w_{i+1},\\
 D_n(v;t_2)&=(I-\widehat T_n)(I+\widehat T_n)^{-1}.
 \end{split}
\end{equation}
Here and below the signs are those of the physical cycle.  Define
$\mathcal J_2(v;t)=v$ and, for $3\le n\le M$, define
\begin{equation}\label{eq:v67-off-fixed-recursion}
 \mathcal J_n(v;t)=D_n(v;t_2)
  \{t_n-\mathcal R_n(\mathcal J_2(v;t),\ldots,
                                      \mathcal J_{n-1}(v;t))\}.
\end{equation}
The lower-jet remainder is the actual smooth function of
Theorem~\ref{thm:v65-cyclic-jets}.  When $v=z(t_2)$, the Schur
identity identifies $\widehat\sigma$ with the actual orbit factors,
so this is exactly the finite-jet inverse.  Away from the fixed point
\eqref{eq:v67-off-fixed-recursion} is an explicitly chosen algebraic
extension of that inverse; no nonstationary tuple is asserted to solve
the physical half-line problem.  This prescription avoids leaving
unspecified how an approximate curvature is used in the later blocks.

Choose a common compact positive curvature box and, inductively, compact
boxes for the finitely many lower jets and action coefficients.  They
can contain all exact reconstructions and their sufficiently small
perturbations: smoothness, compactness and the finite number of steps
suffice.  Enlarge these boxes slightly so the line segments used below
remain in them.  On these boxes let $\alpha<1$ bound
$|\widehat\sigma_i|$, and let $d_*$ bound
$k_i/(k_i+t_{2,i+1}+v_{i+1})^2$.  Then
\begin{equation}\label{eq:v67-block-lipschitz}
 \begin{split}
 \|D_n\|_\infty&\le b_n:=\frac{1+\alpha^n}{1-\alpha^n},\\
 \|D_vD_n\|&\le\beta_n:=
            \frac{2n\alpha^{n-1}d_*}{(1-\alpha^n)^2}.
 \end{split}
\end{equation}
Indeed $D_n=2(I+\widehat T_n)^{-1}-I$, and differentiating this
identity places $D_v\widehat T_n$ between two resolvents.  Its norm
is at most $n\alpha^{n-1}d_*$.  Thus no commutation of a matrix
variation with the cyclic shift has been assumed.

Let $U_n$ bound $\|t_n-\mathcal R_n(w_{<n})\|_\infty$ on the
chosen boxes and let $r_n$ bound the derivative norm of
$\mathcal R_n$ in the joint maximum norm of its lower jets.  Both are
finite by the actual smooth-jet factorization.  The numbers
\begin{equation}\label{eq:v67-propagation-constants}
 L_2=1,\qquad
 L_n=\max\{L_{n-1},\ \beta_nU_n+b_n r_nL_{n-1}\}
                    \quad(3\le n\le M)
\end{equation}
are explicit sufficient propagation constants once these finite smooth
bounds are specified.  They are not claimed to be known from the noisy
law alone.  Subtract \eqref{eq:v67-off-fixed-recursion} at two
curvatures $v,w$, keeping $t$ fixed.  The change in its block is at
most $\beta_nU_n\|v-w\|_\infty$; the change in the remainder
is at most $b_n r_n$ times the maximum lower-jet difference.
Induction therefore proves
\[
 \|\mathcal J_{\le n}(v;t)-\mathcal J_{\le n}(w;t)\|_\infty
                  \le L_n\|v-w\|_\infty.
\]
Apply this with $w=z(t_2)$ and use
\eqref{eq:v67-residual-error}.  The exact-fixed-point reconstruction
already differs from the true jets by
$C_M(\tau^N+\varepsilon)$.  Their sum gives
\eqref{eq:v67-complete-stopping-error}.  Taking the combined error
small keeps all the iterates used here in the common boxes and in
particular keeps $v$ positive.  Alternatively the displayed residual
positivity test certifies that $z(t_2)$ lies in the quadratic image.
For any strictly positive exact fixed point the iterates eventually
pass that test, because their residuals tend to zero.

Error in an observed density is the input here, not a number of
full-phase preparations; the preparation theorems retain their own
observation and calibration hypotheses.  All constants in this proof
are at fixed order.  In particular they do not give an order-uniform
inverse in the unweighted derivative norm.
\end{proof}
